\appendix{Фрагменты исходного кода программы}

В приложении приведены основные фрагменты исходного кода разработанной программной системы. Листинги выбраны таким образом, чтобы показать реализацию компонентной модели, жизненного цикла приложения и сцены, графической подсистемы, подсистемы ввода, скриптов, физической подсистемы, демонстрационного использования и тестирования.
\newpage
\noindent\textbf{Листинг Б.1 -- базовый класс компонента Component.java}
\begin{lstlisting}[language=Java, frame=single]
package qchromatic.jecse.core;

public abstract class Component {
	protected Entity entity;
	private boolean _started;
	private boolean _destroyed;

	public Entity entity () { return entity; }

	public boolean started () { return _started; }

	public boolean destroyed () { return _destroyed; }

	public void onAttach () { }

	public void onDetach () { }

	public void onStart () { }

	public void onStop () { }

	public void onDestroy () { }

	final void attach (Entity entity) {
		if (this.entity == entity)
			return;

		if (this.entity != null && this.entity != entity)
			throw new IllegalStateException("Component is already attached to another entity");

		this.entity = entity;
		_destroyed = false;
		onAttach();
	}

	final void detach (Entity entity) {
		if (this.entity == entity) {
			stopInternal();
			onDetach();
			this.entity = null;
		}
	}

	final void startInternal () {
		if (_started || _destroyed) return;

		_started = true;
		onStart();
	}

	final void stopInternal () {
		if (!_started) return;

		_started = false;
		onStop();
	}

	final void destroyInternal () {
		if (_destroyed) return;

		stopInternal();
		_destroyed = true;
		onDestroy();

		if (this instanceof Disposable disposable)
			disposable.destroy();
	}
}
\end{lstlisting}

\newpage
\noindent\textbf{Листинг Б.2 -- базовый класс системы обработки System.java}
\begin{lstlisting}[language=Java, frame=single]
package qchromatic.jecse.core;

import qchromatic.jecse.engine.Scene;

public abstract class System {
	protected Scene scene;
	private int _order;

	public System () { this(0); }

	protected System (int order) {
		_order = order;
	}

	public void init () { }
	public void start () { }
	public void loop (float dtime) { }
	public void stop () { }

	public void destroy () { }

	public void onEntityAdded (Entity entity) { }
	public void onEntityRemoved (Entity entity) { }

	public void onComponentAdded (Entity entity, Component component) { }
	public void onComponentRemoved (Entity entity, Component component) { }

	public int order () { return _order; }

	public System order (int order) {
		_order = order;
		return this;
	}

	public System scene (Scene scene) {
		if (scene == null) return this;

		this.scene = scene;
		return this;
	}
}
\end{lstlisting}

\newpage
\noindent\textbf{Листинг Б.3 -- сущность и управление компонентами Entity.java}
\begin{lstlisting}[language=Java, frame=single]
package qchromatic.jecse.core;

import java.util.*;
import java.util.concurrent.atomic.AtomicLong;

public class Entity {
	private static final AtomicLong NEXT_ID = new AtomicLong(1);

	private final String _id;
	private final Map<Class<? extends Component>, Component> _components;

	private EntityObserver _observer;

	public Entity () { this("entity-" + NEXT_ID.getAndIncrement()); }
	public Entity (String id) {
		if (id == null || id.isBlank())
			throw new IllegalArgumentException("Entity id cannot be null or blank");

		_id = id;
		_components = new LinkedHashMap<>();
	}

	public String id () { return _id; }

	public void observer (EntityObserver observer) { _observer = observer; }

	public <T extends Component> Entity addComponent (T component) {
		if (component == null)
			throw new IllegalArgumentException("Component cannot be null");

		if (component.entity != null && component.entity != this)
			throw new IllegalStateException("Component is already attached to another entity");

		if (_components.containsKey(component.getClass()))
			throw new IllegalStateException("Entity already has component: " + component.getClass().getName());

		component.attach(this);
		_components.put(component.getClass(), component);

		if (_observer != null)
			_observer.onComponentAdded(this, component);

		return this;
	}

	public Component[] getAllComponents () {
		return _components.values().toArray(new Component[0]);
	}

	public Collection<Component> components () {
		return Collections.unmodifiableCollection(_components.values());
	}

	public void detachAllComponents () {
		for (Component component : _components.values())
			component.detach(this);
	}

	public void attachAllComponents () {
		for (Component component : _components.values())
			component.attach(this);
	}

	public void startComponents () {
		for (Component component : _components.values())
			component.startInternal();
	}

	public void stopComponents () {
		for (Component component : _components.values())
			component.stopInternal();
	}

	public void destroyComponents () {
		for (Component component : _components.values())
			component.destroyInternal();
	}

	public void startComponent (Component component) {
		if (component != null && _components.containsValue(component))
			component.startInternal();
	}

	public <T extends Component> T getComponent (Class<T> componentClass) {
		if (componentClass == null) return null;
		return componentClass.cast(_components.get(componentClass));
	}

	public <T extends Component> List<T> getComponentsInheritance(Class<T> componentClass) {
		List<T> result = new ArrayList<>();
		for (Component component : _components.values()) {
			if (componentClass.isAssignableFrom(component.getClass())) {
				result.add(componentClass.cast(component));
			}
		}

		return result;
	}

	public boolean hasComponent (Class<? extends Component> componentClass) {
		if (componentClass == null) return false;
		return _components.containsKey(componentClass);
	}

	public boolean hasComponentInheritance(Class<? extends Component> componentClass) {
		if (componentClass == null) return false;

		for (Component component : _components.values()) {
			if (componentClass.isAssignableFrom(component.getClass())) {
				return true;
			}
		}
		return false;
	}

	public void removeComponent (Class<? extends Component> componentClass) {
		if (componentClass == null) return;

		Component component = _components.remove(componentClass);

		if (component == null) return;

		if (_observer != null)
			_observer.onComponentRemoved(this, component);
		else
			finishComponentRemoval(component);
	}

	public void finishComponentRemoval (Component component) {
		if (component == null) return;
		component.destroyInternal();
		component.detach(this);
	}
}
\end{lstlisting}

\newpage
\noindent\textbf{Листинг Б.4 -- фрагменты управления сценой Scene.java}
\begin{lstlisting}[language=Java, frame=single]
public final class Scene {
    	public void init () {
    		ensureNotDisposed();
    		initializeIfNeeded();
    		start();
    	}

    	public void start () {
    		ensureNotDisposed();
    		initializeIfNeeded();

    		if (_running) return;

    		_running = true;
    		for (Entity entity : entitySnapshot())
    			entity.startComponents();

    		for (System system : systemSnapshot())
    			system.start();

    		flushStructuralChanges();
    	}

    	public void loop (float dtime) {
    		if (!_running) return;

    		flushStructuralChanges();

    		_updating = true;
    		try {
    			for (System system : systemSnapshot())
    				system.loop(dtime);
    		} finally {
    			_updating = false;
    		}

    		flushStructuralChanges();
    	}

    	public void stop () {
    		if (!_running) return;

    		flushStructuralChanges();
    		_running = false;

    		for (System system : systemSnapshot())
    			system.stop();

    		for (Entity entity : entitySnapshot())
    			entity.stopComponents();
    	}

    	public void dispose () {
    		if (_disposed) return;

    		stop();

    		for (System system : systemSnapshot())
    			system.destroy();

    		for (Entity entity : entitySnapshot())
    			disposeEntity(entity);

    		_pendingStructuralChanges.clear();
    		_componentIndex.clear();
    		_queries.clear();
    		_systems.clear();
    		_systemMap.clear();
    		_entities.clear();

    		_initialized = false;
    		_disposed = true;
    	}

    	public void addEntity (Entity entity) {
    		ensureNotDisposed();
    		if (entity == null)
    			throw new IllegalArgumentException("Entity cannot be null");

    		if (_updating || _flushingStructuralChanges) {
    			_pendingStructuralChanges.add(() -> addEntityNow(entity));
    			return;
    		}

    		addEntityNow(entity);
    	}

    	public void removeEntity (String id) {
    		ensureNotDisposed();

    		if (_updating || _flushingStructuralChanges) {
    			_pendingStructuralChanges.add(() -> removeEntityNow(id));
    			return;
    		}

    		removeEntityNow(id);
    	}

    	public void addSystem (System system) {
    		ensureNotDisposed();
    		if (system == null)
    			throw new IllegalArgumentException("System cannot be null");

    		if (_updating || _flushingStructuralChanges) {
    			_pendingStructuralChanges.add(() -> addSystemNow(system));
    			return;
    		}

    		addSystemNow(system);
    	}

    	public void removeSystem (Class<? extends System> systemClass) {
    		ensureNotDisposed();

    		if (_updating || _flushingStructuralChanges) {
    			_pendingStructuralChanges.add(() -> removeSystemNow(systemClass));
    			return;
    		}

    		removeSystemNow(systemClass);
    	}

    	public final EntityQuery query (Class<? extends Component>... componentClasses) {
    		return createQuery(false, Arrays.asList(componentClasses));
    	}

    	public final EntityQuery queryInheritance (Class<? extends Component>... componentClasses) {
    		return createQuery(true, Arrays.asList(componentClasses));
    	}

    	private void flushStructuralChanges () {
    		if (_flushingStructuralChanges || _pendingStructuralChanges.isEmpty()) return;

    		_flushingStructuralChanges = true;
    		try {
    			while (!_pendingStructuralChanges.isEmpty()) {
    				List<Runnable> changes = List.copyOf(_pendingStructuralChanges);
    				_pendingStructuralChanges.clear();

    				for (Runnable change : changes)
    					change.run();
    			}
    		} finally {
    			_flushingStructuralChanges = false;
    		}
    	}
}
\end{lstlisting}

\newpage
\noindent\textbf{Листинг Б.5 -- главный цикл приложения Application.java}
\begin{lstlisting}[language=Java, frame=single]
	public void run (Scene startScene) {
		_context.use();

		try {
			if (startScene != null)
				_context.sceneManager().load(startScene);
			else if (!_context.sceneManager().hasActiveScene())
				throw new RuntimeException("No start scene");

			updateCameraAspect();
			_window.show();

			_context.time().reset((float) glfwGetTime());
			while (!_window.shouldClose()) {
				_window.clear(_context.config().clearColor());

				float currentTime = (float) glfwGetTime();
				_context.time().update(currentTime);
				_context.sceneManager().activeScene().loop(_context.time().deltaTime());
				_context.sceneManager().applyPendingChanges();
				updateCameraAspect();
				_context.input().update();
				_window.update();
			}
		} finally {
			_context.destroy();
		}
	}
\end{lstlisting}

\newpage
\noindent\textbf{Листинг Б.6 -- управление сменой сцен SceneManager.java}
\begin{lstlisting}[language=Java, frame=single]
package qchromatic.jecse.engine;

public final class SceneManager {
	private static final SceneManager DEFAULT = new SceneManager();

	private Scene _activeScene;
	private Scene _pendingScene;
	private boolean _pendingReplace;

	public Scene activeScene () {
		if (_activeScene == null) throw new RuntimeException("No active scene");
		return _activeScene;
	}

	public boolean hasActiveScene () {
		return _activeScene != null;
	}

	public void load (Scene scene) {
		if (scene == null) return;
		clearPendingChange();

		if (_activeScene == scene) {
			_activeScene.start();
			return;
		}

		if (_activeScene != null)
			_activeScene.stop();

		_activeScene = scene;
		_activeScene.start();
	}

	public void replace (Scene scene) {
		if (scene == null) return;
		clearPendingChange();

		if (_activeScene == scene) {
			_activeScene.start();
			return;
		}

		if (_activeScene != null) {
			Scene previousScene = _activeScene;
			_activeScene = null;
			previousScene.dispose();
		}

		_activeScene = scene;
		_activeScene.start();
	}

	public void requestLoad (Scene scene) {
		if (scene == null) return;

		_pendingScene = scene;
		_pendingReplace = false;
	}

	public void requestReplace (Scene scene) {
		if (scene == null) return;

		_pendingScene = scene;
		_pendingReplace = true;
	}

	public void applyPendingChanges () {
		if (_pendingScene == null) return;

		Scene scene = _pendingScene;
		boolean replace = _pendingReplace;
		clearPendingChange();

		if (replace)
			replace(scene);
		else
			load(scene);
	}

	public void unload () {
		if (_activeScene == null) return;
		clearPendingChange();

		_activeScene.stop();
		_activeScene = null;
	}

	public void disposeActiveScene () {
		if (_activeScene == null) return;
		clearPendingChange();

		Scene scene = _activeScene;
		_activeScene = null;
		scene.dispose();
	}

	public static Scene getActiveScene () {
		return currentManager().activeScene();
	}

	public static boolean hasScene () {
		return currentManager().hasActiveScene();
	}

	public static void loadScene (Scene scene) {
		currentManager().load(scene);
	}

	public static void replaceScene (Scene scene) {
		currentManager().replace(scene);
	}

	public static void requestLoadScene (Scene scene) {
		currentManager().requestLoad(scene);
	}

	public static void requestReplaceScene (Scene scene) {
		currentManager().requestReplace(scene);
	}

	public static void unloadScene () {
		currentManager().unload();
	}

	public static void disposeScene () {
		currentManager().disposeActiveScene();
	}

	private static SceneManager currentManager () {
		EngineContext context = EngineContext.currentOrNull();
		return context == null ? DEFAULT : context.sceneManager();
	}

	private void clearPendingChange () {
		_pendingScene = null;
		_pendingReplace = false;
	}
}
\end{lstlisting}

\newpage
\noindent\textbf{Листинг Б.7 -- фрагменты графической подсистемы RenderSystem.java}
\begin{lstlisting}[language=Java, frame=single]
public final class RenderSystem {
    	public void init () {
    		_renderables = scene.query(Transform.class, MeshRenderer.class);
    		_cameras = scene.query(Transform.class, Camera.class);
    		_lights = scene.query(DirectionalLight.class);
    	}

    	public void loop (float dtime) {
    		Entity cameraEntity = selectCamera();
    		if (cameraEntity == null) return;

    		Camera camera = cameraEntity.getComponent(Camera.class);
    		LightState light = selectLight();

    		buildRenderQueue(camera);
    		sortRenderQueue();

    		for (RenderCommand command : _renderQueue)
    			renderCommand(command, camera, light);

    		checkOpenGLErrors();
    	}

    	private Entity selectCamera () {
    		Entity selected = null;
    		int selectedPriority = Integer.MIN_VALUE;

    		for (Entity entity : _cameras) {
    			Camera camera = entity.getComponent(Camera.class);
    			if (camera == null || !camera.enabled()) continue;

    			if (selected == null || camera.priority() > selectedPriority) {
    				selected = entity;
    				selectedPriority = camera.priority();
    			}
    		}

    		return selected;
    	}

    	private LightState selectLight () {
    		for (Entity entity : _lights) {
    			DirectionalLight light = entity.getComponent(DirectionalLight.class);
    			if (light != null && light.enabled())
    				return new LightState(light.direction(), light.color().multiplied(light.intensity()), true);
    		}

    		return new LightState(new Vec3(0f, -1f, 0f), new Vec4(1f), false);
    	}

    	private void buildRenderQueue (Camera camera) {
    		_renderQueue.clear();

    		for (Entity entity : _renderables) {
    			MeshRenderer meshRenderer = entity.getComponent(MeshRenderer.class);
    			Transform transform = entity.getComponent(Transform.class);

    			if (meshRenderer == null || transform == null || !meshRenderer.enabled()) continue;
    			if ((camera.cullingMask() & meshRenderer.layer()) == 0) continue;

    			Mesh mesh = meshRenderer.mesh();
    			Material material = meshRenderer.material();

    			if (mesh == null || material == null) continue;
    			if (!isVisible(camera, transform, meshRenderer)) continue;

    			_renderQueue.add(new RenderCommand(mesh, material, transform, meshRenderer.renderQueue()));
    		}
    	}

    	private void sortRenderQueue () {
    		_renderQueue.sort(Comparator
    				.comparingInt((RenderCommand command) -> command.renderQueue)
    				.thenComparingInt(command -> command.material.shader().handler())
    				.thenComparingInt(command -> command.material.texture().id())
    				.thenComparingInt(command -> java.lang.System.identityHashCode(command.mesh)));
    	}

    	private void renderCommand (RenderCommand command, Camera camera, LightState light) {
    		MeshGpuResource resource = getMeshGpuResource(command.mesh);
    		Material material = command.material;

    		material.use();

    		Shader shader = material.shader();
    		shader.setUniform("u_model", command.transform.getModelMatrix().transposed());
    		shader.setUniform("u_view", camera.getViewMatrix().transposed());
    		shader.setUniform("u_projection", camera.getProjectionMatrix().transposed());
    		shader.setUniform("u_lightDirection", light.direction);
    		shader.setUniform("u_lightColor", light.color);
    		shader.setUniform("u_ambientColor", _ambientColor);
    		shader.setUniform("u_lit", light.enabled);

    		glBindVertexArray(resource.vaoId());
    		glDrawElements(GL_TRIANGLES, resource.indexCount(), GL_UNSIGNED_INT, 0);
    		glBindVertexArray(0);
    	}
}
\end{lstlisting}

\newpage
\noindent\textbf{Листинг Б.8 -- фрагменты физической подсистемы PhysicsSystem.java}
\begin{lstlisting}[language=Java, frame=single]
public final class PhysicsSystem {
    	public void loop (float dtime) {
    		if (dtime <= 0f) return;

    		_accumulator += Math.min(dtime, _fixedTimeStep * _maxSubSteps);

    		int steps = 0;
    		while (_accumulator >= _fixedTimeStep && steps < _maxSubSteps) {
    			step(_fixedTimeStep);
    			_accumulator -= _fixedTimeStep;
    			steps++;
    		}

    		if (steps == _maxSubSteps && _accumulator >= _fixedTimeStep)
    			_accumulator = 0f;

    		if (_debugDrawColliders)
    			drawColliderBounds();
    	}

    	private void step (float dtime) {
    		integrateForces(dtime);
    		integrateTransforms(dtime);
    		solveVelocity();
    		solvePosition();
    		updateSleeping(dtime);
    	}

    	private void integrateForces (float dtime) {
    		for (Entity entity : _bodies) {
    			Transform transform = entity.getComponent(Transform.class);
    			Rigidbody body = entity.getComponent(Rigidbody.class);
    			if (transform == null || body == null || !body.enabled()) continue;

    			if (body.staticBody()) {
    				body.velocity(0f, 0f, 0f);
    				body.angularVelocity(0f, 0f, 0f);
    				body.clearForces();
    				continue;
    			}

    			if (body.sleeping()) {
    				body.clearForces();
    				continue;
    			}

    			if (body.dynamic()) {
    				Vec3 velocity = body.velocity();
    				Vec3 acceleration = body.force().multiplied(body.inverseMass());
    				if (body.useGravity())
    					acceleration.add(_gravity.multiplied(body.gravityScale()));

    				velocity.add(acceleration.multiplied(dtime));
    				velocity.mul(Math.max(0f, 1f - body.linearDamping() * dtime));
    				body.velocity(velocity);

    				Vec3 angularVelocity = body.angularVelocity();
    				if (!body.lockRotation()) {
    					Vec3 angularAcceleration = applyInverseInertia(entity, transform, body, body.torque());
    					angularVelocity.add(angularAcceleration.multiplied(dtime));
    					angularVelocity.mul(Math.max(0f, 1f - body.angularDamping() * dtime));
    					body.angularVelocity(angularVelocity);
    				} else {
    					body.angularVelocity(0f, 0f, 0f);
    				}
    			}

    			clampBodySpeeds(body);
    			body.clearForces();
    		}
    	}

    	private void integrateTransforms (float dtime) {
    		for (Entity entity : _bodies) {
    			Transform transform = entity.getComponent(Transform.class);
    			Rigidbody body = entity.getComponent(Rigidbody.class);
    			if (transform == null || body == null || !body.enabled() || body.staticBody() || body.sleeping()) continue;

    			transform.position(transform.position().added(body.velocity().multiplied(dtime)));

    			Vec3 angularVelocity = body.angularVelocity();
    			if (!body.lockRotation() && angularVelocity.lengthSquared() > 0f) {
    				float angleRad = angularVelocity.length() * dtime;
    				Vec3 axis = angularVelocity.normalized();
    				Quaternion delta = Quaternion.axisAngle(axis, (float) Math.toDegrees(angleRad));
    				transform.rotation(delta.multiplied(transform.rotation()).normalized());
    			}
    		}
    	}

    	private List<Contact> contacts () {
    		List<ColliderEntry> entries = colliderEntries();
    		List<Contact> result = new ArrayList<>();

    		for (int i = 0; i < entries.size(); i++) {
    			for (int j = i + 1; j < entries.size(); j++) {
    				ColliderEntry a = entries.get(i);
    				ColliderEntry b = entries.get(j);

    				if (!canCollide(a.collider, b.collider)) continue;

    				Contact contact = contact(a, b);
    				if (contact != null)
    					result.add(contact);
    			}
    		}

    		return result;
    	}

    	private Contact contact (ColliderEntry a, ColliderEntry b) {
    		if (!aabbOverlap(a.shape, b.shape)) return null;

    		SatResult result = new SatResult();
    		Vec3 centerDelta = subtract(b.shape.center, a.shape.center);

    		for (int i = 0; i < 3; i++) {
    			if (!testAxis(a.shape.axes[i], centerDelta, a.shape, b.shape, result)) return null;
    			if (!testAxis(b.shape.axes[i], centerDelta, a.shape, b.shape, result)) return null;
    		}

    		for (int i = 0; i < 3; i++) {
    			for (int j = 0; j < 3; j++) {
    				Vec3 axis = a.shape.axes[i].crossed(b.shape.axes[j]);
    				if (axis.lengthSquared() > EPSILON && !testAxis(axis, centerDelta, a.shape, b.shape, result))
    					return null;
    			}
    		}

    		if (result.normal.lengthSquared() <= 0f)
    			return null;

    		return new Contact(a, b, result.normal, result.penetration, contactPoints(a.shape, b.shape, result.normal, result.penetration));
    	}

    	private boolean testAxis (Vec3 axis, Vec3 centerDelta, BoxShape a, BoxShape b, SatResult result) {
    		Vec3 normal = axis.normalized();
    		float radiusA = projectionRadius(a, normal);
    		float radiusB = projectionRadius(b, normal);
    		float distance = centerDelta.dot(normal);
    		float overlap = radiusA + radiusB - Math.abs(distance);

    		if (overlap <= 0f)
    			return false;

    		if (overlap < result.penetration) {
    			result.penetration = overlap;
    			result.normal = distance >= 0f ? normal : normal.multiplied(-1f);
    		}

    		return true;
    	}

    	private void resolveVelocityAtPoint (Contact contact, Vec3 point) {
    		float invMassA = inverseMass(contact.a.body);
    		float invMassB = inverseMass(contact.b.body);
    		Vec3 centerA = worldCenterOfMass(contact.a);
    		Vec3 centerB = worldCenterOfMass(contact.b);
    		Vec3 rA = subtract(point, centerA);
    		Vec3 rB = subtract(point, centerB);

    		Vec3 relativeVelocity = subtract(pointVelocity(contact.b, rB), pointVelocity(contact.a, rA));
    		float velocityAlongNormal = relativeVelocity.dot(contact.normal);
    		float velocityBias = contactVelocityBias(contact, velocityAlongNormal);
    		if (velocityAlongNormal >= velocityBias)
    			return;

    		float denominator = invMassA + invMassB
    				+ angularDenominator(contact.a, rA, contact.normal)
    				+ angularDenominator(contact.b, rB, contact.normal);
    		if (denominator <= 0f)
    			return;

    		float impulseMagnitude = (velocityBias - velocityAlongNormal) / denominator;
    		applyImpulse(contact.a, contact.normal.multiplied(-impulseMagnitude), point);
    		applyImpulse(contact.b, contact.normal.multiplied(impulseMagnitude), point);

    		relativeVelocity = subtract(pointVelocity(contact.b, rB), pointVelocity(contact.a, rA));
    		Vec3 tangent = subtract(relativeVelocity, contact.normal.multiplied(relativeVelocity.dot(contact.normal)));
    		if (tangent.lengthSquared() > EPSILON) {
    			tangent.normalize();
    			float tangentDenominator = invMassA + invMassB
    					+ angularDenominator(contact.a, rA, tangent)
    					+ angularDenominator(contact.b, rB, tangent);
    			if (tangentDenominator > 0f) {
    				float tangentImpulseMagnitude = -relativeVelocity.dot(tangent) / tangentDenominator;
    				float friction = (float) Math.sqrt(contact.a.collider.friction() * contact.b.collider.friction());
    				float maxFriction = Math.abs(impulseMagnitude) * friction;
    				tangentImpulseMagnitude = clamp(tangentImpulseMagnitude, -maxFriction, maxFriction);

    				applyImpulse(contact.a, tangent.multiplied(-tangentImpulseMagnitude), point);
    				applyImpulse(contact.b, tangent.multiplied(tangentImpulseMagnitude), point);
    			}
    		}
    	}

    	private void resolvePosition (Contact contact) {
    		float invMassA = inverseMass(contact.a.body);
    		float invMassB = inverseMass(contact.b.body);
    		float invMassSum = invMassA + invMassB;
    		if (invMassSum <= 0f) return;

    		float percent = 0.2f;
    		float correction = Math.min(0.05f, Math.max(0f, contact.penetration - POSITION_SLOP) * percent);
    		if (correction <= 0f) return;

    		if (invMassA > 0f)
    			move(contact.a.transform, contact.normal.multiplied(-correction * invMassA / invMassSum));
    		if (invMassB > 0f)
    			move(contact.b.transform, contact.normal.multiplied(correction * invMassB / invMassSum));
    	}

    	private void updateSleeping (float dtime) {
    		Map<Rigidbody, Boolean> supported = new IdentityHashMap<>();
    		Map<Rigidbody, Boolean> touchedByActiveBody = new IdentityHashMap<>();
    		for (Contact contact : contacts()) {
    			if (contact.a.collider.trigger() || contact.b.collider.trigger()) continue;
    			if (isSupported(contact, contact.a, contact.normal.multiplied(-1f)) && isStableSupport(contact.b.body))
    				supported.put(contact.a.body, true);
    			if (isSupported(contact, contact.b, contact.normal) && isStableSupport(contact.a.body))
    				supported.put(contact.b.body, true);
    			if (isActiveBody(contact.b.body) && contact.a.body != null)
    				touchedByActiveBody.put(contact.a.body, true);
    			if (isActiveBody(contact.a.body) && contact.b.body != null)
    				touchedByActiveBody.put(contact.b.body, true);
    		}

    		for (Entity entity : _bodies) {
    			Rigidbody body = entity.getComponent(Rigidbody.class);
    			if (body == null || !body.enabled() || !body.dynamic()) continue;

    			if (!supported.containsKey(body)) {
    				_sleepTimers.remove(body);
    				if (body.sleeping())
    					body.wakeUp();
    				continue;
    			}

    			if (touchedByActiveBody.containsKey(body)) {
    				_sleepTimers.remove(body);
    				if (body.sleeping())
    					body.wakeUp();
    				continue;
    			}

    			boolean slow = body.velocity().lengthSquared() <= _sleepLinearThreshold * _sleepLinearThreshold
    					&& body.angularVelocity().lengthSquared() <= _sleepAngularThreshold * _sleepAngularThreshold;
    			if (!slow) {
    				_sleepTimers.remove(body);
    				continue;
    			}

    			float sleepTime = _sleepTimers.getOrDefault(body, 0f) + dtime;
    			_sleepTimers.put(body, sleepTime);
    			if (sleepTime >= _sleepDelay)
    				body.sleep();
    		}
    	}
}
\end{lstlisting}

\newpage
\noindent\textbf{Листинг Б.9 -- выполнение пользовательских скриптов ScriptSystem.java}
\begin{lstlisting}[language=Java, frame=single]
package qchromatic.jecse.system;

import qchromatic.jecse.component.Script;
import qchromatic.jecse.core.Component;
import qchromatic.jecse.core.Entity;
import qchromatic.jecse.core.EntityQuery;
import qchromatic.jecse.core.System;

import java.util.*;

public class ScriptSystem extends System {
	private final List<Script> _scripts;
	private final Set<Script> _initializedScripts;
	private final Set<Script> _startedScripts;
	private final Map<Entity, List<Script>> _scriptsByEntity;
	private EntityQuery _scriptEntities;
	private boolean _running;

	public ScriptSystem () {
		super(0);
		_scripts = new ArrayList<>();
		_initializedScripts = new LinkedHashSet<>();
		_startedScripts = new LinkedHashSet<>();
		_scriptsByEntity = new LinkedHashMap<>();
		_running = false;
	}

	@Override
	public void init () {
		_scriptEntities = scene.queryInheritance(Script.class);
		_scripts.clear();
		_initializedScripts.clear();
		_startedScripts.clear();
		_scriptsByEntity.clear();

		for (Entity entity : _scriptEntities)
			registerEntity(entity);
	}

	@Override
	public void start () {
		_running = true;
		for (Script script : List.copyOf(_scripts))
			startScript(script);
	}

	@Override
	public void loop (float dtime) {
		for (Script script : List.copyOf(_scripts))
			script.loop(dtime);
	}

	@Override
	public void stop () {
		_running = false;
		for (Script script : List.copyOf(_startedScripts))
			stopScript(script);
	}

	@Override
	public void destroy () {
		for (Script script : List.copyOf(_scripts))
			destroyScript(script);

		_scripts.clear();
		_initializedScripts.clear();
		_startedScripts.clear();
		_scriptsByEntity.clear();
	}

	@Override
	public void onEntityAdded (Entity entity) {
		registerEntity(entity);
	}

	@Override
	public void onEntityRemoved (Entity entity) {
		unregisterEntity(entity, true);
	}

	@Override
	public void onComponentAdded(Entity entity, Component component) {
		if (component instanceof Script script)
			registerScript(entity, script);
	}

	@Override
	public void onComponentRemoved(Entity entity, Component component) {
		if (component instanceof Script script)
			unregisterScript(entity, script, true);
	}

	private void registerEntity (Entity entity) {
		for (Script script : entity.getComponentsInheritance(Script.class))
			registerScript(entity, script);
	}

	private void registerScript (Entity entity, Script script) {
		List<Script> entityScripts = _scriptsByEntity.computeIfAbsent(entity, ignored -> new ArrayList<>());
		if (entityScripts.contains(script)) return;

		entityScripts.add(script);
		_scripts.add(script);
		initScript(script);

		if (_running)
			startScript(script);
	}

	private void unregisterEntity (Entity entity, boolean destroy) {
		List<Script> entityScripts = _scriptsByEntity.remove(entity);
		if (entityScripts == null) return;

		for (Script script : List.copyOf(entityScripts))
			unregisterScript(entity, script, destroy);
	}

	private void unregisterScript (Entity entity, Script script, boolean destroy) {
		List<Script> entityScripts = _scriptsByEntity.get(entity);
		if (entityScripts != null) {
			entityScripts.remove(script);
			if (entityScripts.isEmpty())
				_scriptsByEntity.remove(entity);
		}

		_scripts.remove(script);

		if (destroy)
			destroyScript(script);
	}

	private void initScript (Script script) {
		if (_initializedScripts.add(script))
			script.init();
	}

	private void startScript (Script script) {
		initScript(script);
		if (_startedScripts.add(script))
			script.start();
	}

	private void stopScript (Script script) {
		if (_startedScripts.remove(script))
			script.stop();
	}

	private void destroyScript (Script script) {
		stopScript(script);

		if (_initializedScripts.remove(script))
			script.destroy();
	}
}
\end{lstlisting}

\newpage
\noindent\textbf{Листинг Б.10 -- обработка пользовательского ввода InputState.java}
\begin{lstlisting}[language=Java, frame=single]
package qchromatic.jecse.engine;

import qchromatic.jecse.common.Vec2;

public final class InputState {
	private static final int MAX_KEYS = 512;
	private static final int MAX_MOUSE_BUTTONS = 16;

	private final boolean[] _keys;
	private final boolean[] _keysPrev;
	private final boolean[] _mouseButtons;
	private final boolean[] _mouseButtonsPrev;

	private float _mouseX;
	private float _mouseY;
	private float _mouseDX;
	private float _mouseDY;
	private boolean _hasMousePosition;
	private float _scrollX;
	private float _scrollY;

	public InputState () {
		_keys = new boolean[MAX_KEYS];
		_keysPrev = new boolean[MAX_KEYS];
		_mouseButtons = new boolean[MAX_MOUSE_BUTTONS];
		_mouseButtonsPrev = new boolean[MAX_MOUSE_BUTTONS];
	}

	public void update() {
		java.lang.System.arraycopy(_keys, 0, _keysPrev, 0, MAX_KEYS);
		java.lang.System.arraycopy(_mouseButtons, 0, _mouseButtonsPrev, 0, MAX_MOUSE_BUTTONS);
		_mouseDX = 0f;
		_mouseDY = 0f;
		_scrollX = 0f;
		_scrollY = 0f;
	}

	public void reset () {
		for (int i = 0; i < MAX_KEYS; i++) {
			_keys[i] = false;
			_keysPrev[i] = false;
		}

		for (int i = 0; i < MAX_MOUSE_BUTTONS; i++) {
			_mouseButtons[i] = false;
			_mouseButtonsPrev[i] = false;
		}

		_mouseDX = 0f;
		_mouseDY = 0f;
		_mouseX = 0f;
		_mouseY = 0f;
		_hasMousePosition = false;
		_scrollX = 0f;
		_scrollY = 0f;
	}

	public void setKeyState(int key, boolean pressed) {
		if (key >= 0 && key < MAX_KEYS) _keys[key] = pressed;
	}

	public void setMouseButtonState(int button, boolean pressed) {
		if (button >= 0 && button < MAX_MOUSE_BUTTONS) _mouseButtons[button] = pressed;
	}

	public void setMousePosition(float x, float y) {
		if (!_hasMousePosition) {
			_mouseX = x;
			_mouseY = y;
			_mouseDX = 0f;
			_mouseDY = 0f;
			_hasMousePosition = true;
			return;
		}

		_mouseDX += x - _mouseX;
		_mouseDY += y - _mouseY;
		_mouseX = x;
		_mouseY = y;
	}

	public void setScrollOffset(float xoffset, float yoffset) {
		_scrollX += xoffset;
		_scrollY += yoffset;
	}

	public void resetMousePositionTracking () {
		_hasMousePosition = false;
		_mouseDX = 0f;
		_mouseDY = 0f;
	}

	public boolean getKey (int key) {
		return key >= 0 && key < MAX_KEYS && _keys[key];
	}

	public boolean getKeyDown (int key) {
		return key >= 0 && key < MAX_KEYS && _keys[key] && !_keysPrev[key];
	}

	public boolean getKeyUp (int key) {
		return key >= 0 && key < MAX_KEYS && !_keys[key] && _keysPrev[key];
	}

	public boolean getMouseButton (int button) {
		return button >= 0 && button < MAX_MOUSE_BUTTONS && _mouseButtons[button];
	}

	public boolean getMouseButtonDown (int button) {
		return button >= 0 && button < MAX_MOUSE_BUTTONS && _mouseButtons[button] && !_mouseButtonsPrev[button];
	}

	public boolean getMouseButtonUp (int button) {
		return button >= 0 && button < MAX_MOUSE_BUTTONS && !_mouseButtons[button] && _mouseButtonsPrev[button];
	}

	public Vec2 getMousePosition () {
		return new Vec2(_mouseX, _mouseY);
	}

	public Vec2 getMouseDelta () {
		return new Vec2(_mouseDX, _mouseDY);
	}

	public Vec2 getMouseScroll () {
		return new Vec2(_scrollX, _scrollY);
	}
}
\end{lstlisting}

\newpage
\noindent\textbf{Листинг Б.11 -- фрагменты демонстрационной физической сцены PhysicsScene.java}
\begin{lstlisting}[language=Java, frame=single]
public final class PhysicsScene {
    	private static Scene createScene () {
    		Scene scene = new Scene();

    		scene.addSystem(new ScriptSystem());

    		scene.addSystem(new PhysicsSystem()
    				.gravity(0f, -9.81f, 0f)
    				.fixedTimeStep(1f / 60f)
    				.maxSubSteps(5)
    				.velocityIterations(12)
    				.positionIterations(10)
    				.maxSpeeds(16f, 8f)
    				.sleepThresholds(0.08f, 0.12f, 0.5f)
    				.debugDrawColliders(true));

    		scene.addSystem(new RenderSystem()
    				.ambientColor(new Vec4(0.18f, 0.20f, 0.23f, 1f)));

    		scene.addSystem(new DebugRenderSystem());

    		scene.addEntity(new Entity("sun")
    				.addComponent(new DirectionalLight()
    						.direction(-0.4f, -1f, -0.25f)
    						.color(new Vec4(1f, 0.94f, 0.82f, 1f))
    						.intensity(1.2f)));

    		scene.addEntity(new Entity("camera")
    				.addComponent(new Transform()
    						.position(0f, 4f, 9f)
    						.rotation(Quaternion.euler(-22f, 0f, 0f)))
    				.addComponent(new Camera()
    						.fov(60f)
    						.near(0.1f)
    						.far(80f)
    						.cullingMask(LAYER_WORLD | LAYER_PHYSICS)
    						.priority(10))
    				.addComponent(new FreecamController()
    						.speed(7f)
    						.sensitivity(18f)));

    		addFloor(scene);
    		addWall(scene, "left-wall", -4f, 1f, 0f, 0.25f, 1.2f, 4f);
    		addWall(scene, "right-wall", 4f, 1f, 0f, 0.25f, 1.2f, 4f);
    		addWall(scene, "back-wall", 0f, 1f, -4f, 4f, 1.2f, 0.25f);

    		addNarrowSupport(scene);
    		addDynamicBox(scene, "off-center-box", 0.75f, 3.2f, -1.2f, new Vec4(0.25f, 0.68f, 1f, 1f), new Vec3(0f, 0f, 0f));

    		addDynamicBox(scene, "falling-box-a", -1.2f, 4.5f, 0f, new Vec4(0.25f, 0.68f, 1f, 1f), new Vec3(1.6f, 0f, 0f));
    		addDynamicBox(scene, "falling-box-b", 0.1f, 7f, 0.2f, new Vec4(1f, 0.55f, 0.24f, 1f), new Vec3(-0.6f, 0f, 0.2f));
    		addDynamicBox(scene, "falling-box-c", 3f, 7f, 0.2f, new Vec4(1f, 0f, 1f, 1f), new Vec3(-0.6f, 0f, 0.2f));
    		addDynamicBox(scene, "low-bounce-box", 1.4f, 5.8f, -0.5f, new Vec4(0.55f, 1f, 0.32f, 1f), new Vec3(-1.1f, 0f, 0.5f), 0.08f);

    		scene.addEntity(new Entity("kinematic-platform")
    				.addComponent(new Transform()
    						.position(0f, 0.35f, 3.1f)
    						.scale(1.8f, 0.25f, 0.7f))
    				.addComponent(new Rigidbody()
    						.type(PhysicsBodyType.KINEMATIC)
    						.velocity(0.9f, 0f, 0f)
    						.useGravity(false))
    				.addComponent(new BoxCollider()
    						.halfExtents(0.5f, 0.5f, 0.5f)
    						.friction(0.9f))
    				.addComponent(new MeshRenderer()
    						.mesh(Meshes.cube())
    						.material(new Material().color(new Vec4(1f, 0.9f, 0.28f, 1f)))
    						.layer(LAYER_PHYSICS)
    						.boundsRadius(2f))
    				.addComponent(new PlatformBounceScript(-2.4f, 2.4f)));

    		scene.addEntity(new Entity("physics-grid")
    				.addComponent(new PhysicsGridScript()));

    		scene.addEntity(new Entity("scene-controls")
    				.addComponent(new PhysicsSceneControlScript()));

    		return scene;
    	}
		
    	private static void addDynamicBox (Scene scene, String id, float x, float y, float z, Vec4 color, Vec3 velocity) {
    		addDynamicBox(scene, id, x, y, z, color, velocity, 0.1f);
    	}

    	private static final class PhysicsSceneControlScript extends Script {
    		@Override
    		public void loop (float dtime) {
    			if (!Input.getKeyDown(KeyCode.R)) return;

    			EngineContext context = EngineContext.current();
    			context.window().setCursorDisabled();
    			context.sceneManager().requestReplace(createScene());
    		}
    	}
}
\end{lstlisting}

\newpage
\noindent\textbf{Листинг Б.12 -- фрагменты тестирования физической подсистемы PhysicsSystemTest.java}
\begin{lstlisting}[language=Java, frame=single]
public final class PhysicsSystemTest {
    	private static void testOverhangingBoxFallsToFloor () {
    		Scene scene = new Scene();
    		scene.addSystem(new PhysicsSystem()
    				.gravity(0f, -10f, 0f)
    				.fixedTimeStep(1f / 60f)
    				.velocityIterations(12)
    				.positionIterations(10)
    				.sleepThresholds(0.08f, 0.12f, 0.4f));

    		scene.addEntity(new Entity("floor")
    				.addComponent(new Transform()
    						.position(0f, -0.1f, 0f)
    						.scale(8f, 0.2f, 8f))
    				.addComponent(new BoxCollider()
    						.halfExtents(0.5f, 0.5f, 0.5f)
    						.friction(1f)));

    		scene.addEntity(new Entity("narrow-support")
    				.addComponent(new Transform()
    						.position(0f, 0.25f, 0f))
    				.addComponent(new BoxCollider()
    						.halfExtents(0.5f, 0.25f, 2f)
    						.friction(0.9f)));

    		Rigidbody body = new Rigidbody()
    				.mass(1f)
    				.linearDamping(0.04f)
    				.angularDamping(0.08f);
    		Transform transform = new Transform()
    				.position(0.82f, 1f, 0f);
    		scene.addEntity(new Entity("overhanging-box")
    				.addComponent(transform)
    				.addComponent(body)
    				.addComponent(new BoxCollider()
    						.halfExtents(0.5f, 0.5f, 0.5f)
    						.friction(0.9f)));

    		scene.init();
    		for (int i = 0; i < 120; i++)
    			scene.loop(1f / 60f);

    		if (transform.position().y > 0.8f)
    			throw new AssertionError("overhanging box stayed glued to the support too long: position="
    					+ transform.position().x + "," + transform.position().y + "," + transform.position().z
    					+ " angular=" + body.angularVelocity().x + "," + body.angularVelocity().y + "," + body.angularVelocity().z
    					+ " sleeping=" + body.sleeping());

    		scene.dispose();
    	}

    	private static void testWallContactDoesNotGlueFallingBox () {
    		Scene scene = new Scene();
    		scene.addSystem(new PhysicsSystem()
    				.gravity(0f, -10f, 0f)
    				.fixedTimeStep(1f / 60f)
    				.velocityIterations(10)
    				.positionIterations(10)
    				.maxSpeeds(20f, 10f));

    		scene.addEntity(new Entity("wall")
    				.addComponent(new Transform()
    						.position(0f, 2f, 0f)
    						.scale(0.2f, 5f, 4f))
    				.addComponent(new BoxCollider()
    						.halfExtents(0.5f, 0.5f, 0.5f)
    						.friction(1f)));

    		Rigidbody body = new Rigidbody()
    				.mass(1f)
    				.velocity(-0.3f, 0f, 0f)
    				.linearDamping(0.02f)
    				.angularDamping(0.08f);
    		Transform transform = new Transform()
    				.position(0.55f, 3f, 0f);
    		scene.addEntity(new Entity("wall-touching-box")
    				.addComponent(transform)
    				.addComponent(body)
    				.addComponent(new BoxCollider()
    						.halfExtents(0.5f, 0.5f, 0.5f)
    						.friction(1f)));

    		scene.init();
    		for (int i = 0; i < 180; i++)
    			scene.loop(1f / 60f);

    		if (transform.position().y > 1.5f)
    			throw new AssertionError("box was glued to the wall instead of falling: y=" + transform.position().y);

    		scene.dispose();
    	}

    	private static void testYawSpinSettlesOnFloor () {
    		Scene scene = new Scene();
    		scene.addSystem(new PhysicsSystem()
    				.gravity(0f, -10f, 0f)
    				.fixedTimeStep(1f / 60f)
    				.velocityIterations(12)
    				.positionIterations(10)
    				.maxSpeeds(20f, 10f)
    				.sleepThresholds(0.08f, 0.12f, 0.35f));

    		scene.addEntity(new Entity("ground")
    				.addComponent(new Transform()
    						.position(0f, -0.1f, 0f)
    						.scale(8f, 0.2f, 8f))
    				.addComponent(new BoxCollider()
    						.halfExtents(0.5f, 0.5f, 0.5f)
    						.friction(1f)));

    		Rigidbody body = new Rigidbody()
    				.mass(1f)
    				.angularVelocity(0f, 8f, 0f)
    				.linearDamping(0.05f)
    				.angularDamping(0.08f);
    		Transform transform = new Transform()
    				.position(0f, 0.55f, 0f);
    		scene.addEntity(new Entity("yaw-spinning-box")
    				.addComponent(transform)
    				.addComponent(body)
    				.addComponent(new BoxCollider()
    						.halfExtents(0.5f, 0.5f, 0.5f)
    						.friction(1f)));

    		scene.init();
    		for (int i = 0; i < 600; i++)
    			scene.loop(1f / 60f);

    		if (body.angularVelocity().length() > 0.25f)
    			throw new AssertionError("yaw spinning box did not lose angular velocity");
    		if (Math.abs(transform.position().x) > 0.4f || Math.abs(transform.position().z) > 0.4f)
    			throw new AssertionError("yaw spinning box drifted too far");

    		scene.dispose();
    	}

    	private static void expectNear (String label, float actual, float expected, float tolerance) {
    		if (Math.abs(actual - expected) > tolerance)
    			throw new AssertionError(label + ": expected " + expected + " +/- " + tolerance + ", got " + actual);
    	}
}
\end{lstlisting}
\ifВКР{
\newpage
\addcontentsline{toc}{section}{На отдельных листах (CD-RW в прикрепленном конверте)}
\noindent
\begin{tabular}{p{5.8cm}C{4.8cm}C{4.8cm}}
   Автор ВКР & \lhrulefill{\fill} & \fillcenter\Автор \\
            \setarstrut{\footnotesize}
           & \footnotesize{(подпись, дата)} & \\
            \restorearstrut
   Руководитель ВКР & \lhrulefill{\fill} & \fillcenter\Руководитель \\
            \setarstrut{\footnotesize}
           & \footnotesize{(подпись, дата)} & \\
            \restorearstrut
   Нормоконтроль & \lhrulefill{\fill} & \fillcenter\Нормоконтроль \\
            \setarstrut{\footnotesize}
           & \footnotesize{(подпись, дата)} & \\
            \restorearstrut
\end{tabular}
\vskip 2cm
\begin{center}
\textbf{Место для диска}
\end{center}
}\fi
